\chapter*{Annexe A -- Jenkinsfile}
\addcontentsline{toc}{chapter}{Annexe A -- Jenkinsfile}
\label{annexe-jenkinsfile}

Le script ci-dessous présente la structure du pipeline déclaratif Jenkins, correspondant aux sept étapes décrites en \cref{chap-sprint3}. \emph{Ce fichier est un modèle représentatif de la structure du pipeline~; il doit être remplacé par le contenu exact du \texttt{Jenkinsfile} du projet avant soutenance.}

\begin{lstlisting}[language=bash,caption={Structure du pipeline Jenkins (modèle)}]
pipeline {
  agent any

  environment {
    DOCKERHUB_CREDENTIALS = credentials('dockerhub-credentials')
    SONAR_TOKEN            = credentials('sonarqube-token')
  }

  stages {

    stage('1. Checkout') {
      steps {
        checkout scm
      }
    }

    stage('2. Installation des dependances') {
      steps {
        sh 'cd backend && npm install'
        sh 'cd frontend && npm install'
        sh 'cd agent && pip install -r requirements.txt'
      }
    }

    stage('3. Build') {
      steps {
        sh 'cd frontend && npm run build'
      }
    }

    stage('4. Analyse qualite (SonarQube)') {
      steps {
        withSonarQubeEnv('sonarqube-server') {
          sh 'sonar-scanner'
        }
      }
    }

    stage('5. Scan securite (Trivy)') {
      steps {
        sh 'docker build -t mariemchaabani/monitoring-backend:${BUILD_NUMBER} ./backend'
        sh 'trivy image --severity CRITICAL,HIGH --exit-code 1 mariemchaabani/monitoring-backend:${BUILD_NUMBER}'
      }
    }

    stage('6. Build et push des images Docker') {
      steps {
        sh 'docker build -t mariemchaabani/monitoring-frontend:${BUILD_NUMBER} ./frontend'
        sh 'docker login -u $DOCKERHUB_CREDENTIALS_USR -p $DOCKERHUB_CREDENTIALS_PSW'
        sh 'docker push mariemchaabani/monitoring-backend:${BUILD_NUMBER}'
        sh 'docker push mariemchaabani/monitoring-frontend:${BUILD_NUMBER}'
      }
    }

    stage('7. Deploiement (GitOps)') {
      steps {
        sh 'git commit -am "chore: bump image tag to ${BUILD_NUMBER}" && git push'
        // La synchronisation effective du cluster est ensuite prise en
        // charge par Argo CD (voir Chapitre 6).
      }
    }
  }

  post {
    failure {
      echo 'Pipeline en echec -- deploiement annule.'
    }
  }
}
\end{lstlisting}
